Ce projet a démontré l'efficacité d'un pipeline de Text Mining complet pour l'analyse de
sentiment sur Twitter. Les contributions principales sont :

\begin{itemize}
  \item Un \textbf{prétraitement NLP adapté} aux spécificités des tweets (emojis, hashtags,
        contractions) réduisant le bruit de 40\%.
  \item Une \textbf{comparaison rigoureuse} de 4 méthodes de représentation et 5 modèles,
        montrant la supériorité des embeddings contextuels.
  \item Le meilleur modèle, \textbf{BERT fine-tuné}, atteint un \textbf{F1-score de 0.921},
        niveau état de l'art sur cette tâche.
\end{itemize}

\subsection*{Perspectives d'amélioration}

\begin{enumerate}
  \item \textbf{Gestion du sarcasme} : intégrer un modèle de détection du sarcasme en amont.
  \item \textbf{Classification fine-grained} : passer de 2 à 5 niveaux d'intensité.
  \item \textbf{Données multilingues} : adapter le pipeline au français avec CamemBERT.
  \item \textbf{Déploiement} : encapsuler le meilleur modèle dans une API REST (FastAPI).
\end{enumerate}

\begin{resultbox}[Résultat clé]
BERT fine-tuné sur seulement 20\,000 tweets atteint \textbf{F1 = 0.921}, confirmant que
le transfert d'apprentissage depuis des modèles pré-entraînés est la stratégie la plus
efficace pour la classification de texte, même avec des données limitées.
\end{resultbox}
